\section{Artifact and Reproduction Details}

\subsection{Delivered directory}

The package has four top-level evidence surfaces. \texttt{paper/source}
contains the human-editable manuscript, generated value macros, references,
tables, vector-figure sources, compiled figures, and the ACM PDF.
\texttt{frozen-inputs} contains copies of the campaign manifest, both unseen
round results, the systems preregistration, matrix manifest, aggregate result,
contribution gate, and a compact paper-values file.
\texttt{dossier} contains the complete \routea{} and \routeb{} directories,
including all raw metrics and per-cell reports. \texttt{reproduction}
contains the environment lock, a standalone Python script, exact commands,
and the independent reproduction report.

The delivery index links every surface. A separate artifact manifest records
the SHA-256 digest and byte size of every file, including the complete
campaign dossier. The paper-package manifest binds the route lineage,
systems-result hash, systems-gate hash, audit hash, and all output locations.
The manifest and audit are themselves canonical-model hashes. They always
store \texttt{external\_submission\_authorized=false}.

\subsection{Clean-directory procedure}

The paper builder launches the following procedure in a directory that must
be absent or empty:

\begin{enumerate}
  \item load the frozen paired differences and aggregate values;
  \item recompute the paired mean and the exact
  \BootstrapResamples{}-resample percentile bootstrap using the registered
  seed;
  \item compare every recomputed scalar with the frozen value to numerical
  tolerance \(10^{-12}\);
  \item copy the complete paper source rather than using build outputs from
  the primary directory;
  \item compile the five standalone vector figures; and
  \item compile the ACM two-column paper and record its digest.
\end{enumerate}

The process writes a machine-readable report and exits nonzero if any
statistic differs, a figure fails, or the manuscript fails. The primary
package copies that report but does not infer success from process output
alone. The final auditor loads the report and checks its explicit status.

\subsection{Round-level evidence contract}

For each \routea{} round, the dossier includes:

\begin{itemize}
  \item \path{observation.json}, \path{failure_diagnosis.json}, and
  \path{hypothesis.json};
  \item \path{screening.json}, \path{preregistration.json}, and
  \path{frozen_protocol.json};
  \item development and unseen results with all cell paths and hashes;
  \item deterministic contribution-gate and round-decision records;
  \item research, failure, loop, and validation reports;
  \item evidence-map JSON, tables, figures, and manuscript version; and
  \item the round manifest linking its parent and child records.
\end{itemize}

Round 1 changes the derivative estimator and coefficient aggregation relative
to the parent failure. Round 2 changes both the derivative representation and
the cross-output support operation relative to Round 1. The second round does
not reopen the rejected weak-form or support-stability family. Both unseen
panels remain immutable after reveal.

\subsection{Systems cell semantics}

Each systems cell stores task, family, mode, seed, source-evidence hash, policy
evidence, initial and final failure codes, attempted actions, attempt count,
success, recovery, reproduction, claim support, preregistration state, gate
state, memory and feedback use, intervention count, report claim, scientific
result hash, reproduction hash, wall time, cost, and a canonical result hash.

An initial fault changes one Boolean invariant. A planning controller may
repair only a planning-visible fault. A full-loop controller obtains runtime
failure codes. If feedback is available, it attempts the registered repair;
if the repair needs a \vault{} card and memory is absent, it records a blocked
action. The evidence-gate ablation removes \texttt{claim\_supported} from
visible runtime failures but still evaluates that invariant in the final
task-success vector. This accounts for both its nonzero success rate and its
unsupported-claim count.

\subsection{Citation and manuscript audit}

The citation audit parses every \texttt{\textbackslash cite} command and every
BibTeX key, requiring no missing or unused entries and at least 30 references.
Every entry must include author, title, year, and a DOI, URL, or e-print
identifier. The live audit retrieves the registered primary or metadata
source and records status, resolved URL, and content digest. The LaTeX audit
fails on undefined citations or references, failed standalone figures, a
failed paper build, fewer than eight pages, or an overfull box.

The deterministic pre-submission review also checks section completeness,
placeholder text, disallowed promotional phrases, unresolved quantitative
macros, disclosure of both negative \routea{} results, claim-evidence
coverage, and the clean-directory result. Passing this review means only
``ready for human submission review.'' It never means submitted, accepted, or
authorized for release.

\section{Additional Numerical Results}

Table~\ref{tab:all-modes} is generated directly from the frozen
machine-readable aggregate. For clarity, the central values are repeated
here. Full-loop success is \FullLoopSuccess{}, negative recovery is
\FullLoopRecovery{}, and exact reproduction is \FullLoopReproduction{}.
The execute-once success rate is \ExecuteOnceSuccess{}. The paired gain is
\PairedGain{} with interval [\PairedCILower{}, \PairedCIUpper{}].

The ablation success rates are \NoVaultSuccess{} without \vault{} memory,
\NoFeedbackSuccess{} without failure feedback, \NoPreregSuccess{} without
preregistration, and \NoGateSuccess{} without the evidence gate. The last
condition emits \NoGateUnsupported{} unsupported claims. These are descriptive
outcomes of the frozen controlled matrix and should not be read as population
estimates for autonomous research tasks.

\section{Human Review Checklist}

Before any external use, a reviewer should:

\begin{enumerate}
  \item inspect the two negative \routea{} decisions and confirm that the
  systems manuscript does not imply a positive method result;
  \item examine all task-source licenses and the distinction between fresh
  UCI runs and revealed MDBench trace replay;
  \item assess whether the systems contribution is novel and substantial for
  the selected venue;
  \item verify author identity, affiliations, acknowledgments, conflicts, and
  venue disclosure rules;
  \item rerun the standalone reproduction command from a clean directory;
  \item review every claim-evidence edge and all bibliography metadata; and
  \item provide explicit approval before upload, public release, or
  submission.
\end{enumerate}
